\section*{Capítulo \ref{ch:b2:curves} --- Curvas}

\begin{solution}{exo:b2:curves:1}
$\gamma'(t) = (1 - \cos t,\ \sin t)$, logo
\[
\norm{\gamma'(t)}^2 = (1 - \cos t)^2 + \sin^2 t
= 2 - 2\cos t = 4\sin^2\tfrac t2 ,
\]
e $\norm{\gamma'(t)} = 2\sin\frac t2$ (não negativo em $[0,
2\pi]$). Portanto
\[
L = \int_0^{2\pi} 2\sin\tfrac t2\,\dd t
= \Bigl[-4\cos\tfrac t2\Bigr]_0^{2\pi} = 8 :
\]
um arco da cicloide tem \omterm{def:b2:curves:length}{comprimento} $8$ (para uma roda de raio
$1$) --- um resultado célebre de Wren, sem $\pi$ algum à vista.
\end{solution}

\begin{solution}{exo:b2:curves:2}
Com $x = a\cos t$, $y = b\sin t$: $x' = -a\sin t$, $y' = b\cos
t$, $x'' = -a\cos t$, $y'' = -b\sin t$, logo, pela
\cref{prop:b2:curves:kappaformula}
\[
\kappa(t)
= \frac{x'y'' - y'x''}{(x'^2 + y'^2)^{3/2}}
= \frac{ab\sin^2 t + ab\cos^2 t}
       {(a^2\sin^2 t + b^2\cos^2 t)^{3/2}}
= \frac{ab}{(a^2\sin^2 t + b^2\cos^2 t)^{3/2}} .
\]
O denominador é mínimo quando $\sin t = 0$ (valor $b^3$, pontos
$(\pm a, 0)$) e máximo quando $\cos t = 0$ (valor $a^3$, pontos
$(0, \pm b)$), pois $a > b$. Portanto $\kappa$ é máxima nas
extremidades do eixo maior, $\kappa_{\max} = a/b^2$, e mínima nas
extremidades do eixo menor, $\kappa_{\min} = b/a^2$: a elipse
se dobra mais acentuadamente nas pontas de seu eixo longo.
\end{solution}

\begin{solution}{exo:b2:curves:3}
Para o gráfico $\gamma(x) = (x, \cosh x)$: $\norm{\gamma'(x)} =
\sqrt{1 + \sinh^2 x} = \cosh x$, logo o \omterm{def:b2:curves:length}{comprimento de arco} de $0$ a $x$
vale $\int_0^x \cosh u\,\dd u = \sinh x$. \omterm{thm:b2:curves:frenet2d}{Curvatura} de um gráfico
(\cref{prop:b2:curves:kappaformula}):
\[
\kappa(x) = \frac{f''(x)}{(1 + f'(x)^2)^{3/2}}
= \frac{\cosh x}{\cosh^3 x} = \frac{1}{\cosh^2 x} ,
\]
logo $R(x) = \cosh^2 x$, como anunciado. Note a bela coincidência
$R(x) = 1 + s(x)^2$ com $s = \sinh x$ o \omterm{def:b2:curves:length}{comprimento de arco}: o
\omterm{def:b2:curves:curvature}{raio de curvatura} da catenária cresce com o quadrado do \omterm{def:b2:curves:length}{comprimento de
arco} medido a partir do vértice.
\end{solution}

\begin{solution}{exo:b2:curves:4}
$\gamma'(t) = e^t(\cos t - \sin t,\ \sin t + \cos t)$, logo
\[
\langle \gamma(t), \gamma'(t)\rangle = e^{2t}
\bigl(\cos t(\cos t - \sin t) + \sin t(\sin t + \cos t)\bigr)
= e^{2t},
\]
enquanto $\norm{\gamma(t)} = e^t$ e $\norm{\gamma'(t)} =
e^t\sqrt 2$. Portanto
\[
\cos\angle\bigl(\gamma, \gamma'\bigr)
= \frac{e^{2t}}{e^t \cdot e^t\sqrt2} = \frac{1}{\sqrt2} :
\]
a tangente faz sempre o ângulo $\pi/4$ com o raio --- a propriedade
\emph{equiangular} da espiral logarítmica. \omterm{def:b2:curves:length}{Comprimento de arco}
em $(-\infty, 0]$:
\[
\int_{-\infty}^0 \norm{\gamma'(t)}\,\dd t
= \sqrt2\int_{-\infty}^0 e^t\,\dd t = \sqrt 2 ,
\]
finito embora a espiral dê infinitas voltas em torno da
origem. \omterm{thm:b2:curves:frenet2d}{Curvatura}: com $x'y'' - y'x''$ calculado a partir de $\gamma'' =
e^t(-2\sin t,\ 2\cos t)$,
\[
x'y'' - y'x'' = e^{2t}\bigl(2\cos t(\cos t - \sin t)
+ 2\sin t(\sin t + \cos t)\bigr) = 2e^{2t},
\]
logo $\kappa(t) = \dfrac{2e^{2t}}{(e^t\sqrt2)^3} =
\dfrac{1}{e^t\sqrt2}$: a \omterm{thm:b2:curves:frenet2d}{curvatura} vale $1/(\sqrt2\,
\norm{\gamma})$, decaindo à medida que a espiral cresce.
\end{solution}

\begin{solution}{exo:b2:curves:5}
\emph{Primeiro arco:} $\gamma(t) = (t^2,\ t^4 + t^5)$. Derivadas em
$0$: $\gamma'' = (2, 0) \neq 0$, logo $p = 2$. Em seguida $\gamma^{(3)}(0)
= (0, 0)$, $\gamma^{(4)}(0) = (0, 24)$, não colinear com $(2,
0)$: $q = 4$. Ambos pares: \emph{cúspide de segunda espécie} --- os dois
ramos partem na direção $+u = (1,0)$ e ficam do mesmo
lado da tangente. (De fato $y = x^2 \pm x^{5/2}$ nos dois
ramos: mesmo sinal para $x$ pequeno.)

\emph{Segundo arco:} $\gamma(t) = (t^3, t^4)$. $\gamma'(0) =
\gamma''(0) = 0$, $\gamma^{(3)}(0) = (6, 0)$: $p = 3$, ímpar. Em seguida
$\gamma^{(4)}(0) = (0, 24)$: $q = 4$, par. Ímpar--par:
\emph{ponto ordinário} --- apesar da velocidade nula, a
trajetória $y = x^{4/3}$ atravessa a origem suavemente, mantendo-se
acima de sua tangente $y = 0$.
\end{solution}

\begin{solution}{exo:b2:curves:6}
Derive $c(s) = \gamma(s) + \dfrac{1}{\kappa(s)}N(s)$ usando
as fórmulas de Frenet no plano
(\cref{thm:b2:curves:frenet2d}):
\[
c'(s) = T(s) - \frac{\kappa'(s)}{\kappa(s)^2}N(s)
+ \frac{1}{\kappa(s)}\,\bigl(-\kappa(s)T(s)\bigr)
= -\frac{\kappa'(s)}{\kappa(s)^2}\,N(s) ,
\]
os termos tangenciais se cancelando exatamente. Assim a velocidade da
evoluta é carregada por $N(s)$, que dirige a reta normal de
$\gamma$ em $\gamma(s)$ --- e o ponto $c(s)$ está sobre essa mesma
reta normal: a evoluta é a \emph{\omterm{pb:b2:curves:1}{envoltória} das normais}.
(Onde $\kappa' = 0$ a evoluta tem um ponto singular; é isso
que produz as cúspides da evoluta de uma elipse.)
\end{solution}

\begin{solution}{exo:b2:curves:7}
Seja $G(s) = F(s)F(s)^{\mathsf T}$. Então, usando $F' = FA$ e
$(F^{\mathsf T})' = (F')^{\mathsf T} = A^{\mathsf T}F^{\mathsf T}$,
\[
G' = F'F^{\mathsf T} + F(F^{\mathsf T})'
= FAF^{\mathsf T} + FA^{\mathsf T}F^{\mathsf T}
= F(A + A^{\mathsf T})F^{\mathsf T} = 0
\]
por antissimetria. Logo $G$ é constante no intervalo, igual a
$G(s_0) = F(s_0)F(s_0)^{\mathsf T} = I$: $F(s)$ é ortogonal para
todo $s$. (Alternativamente, sem calcular $G'$ até zero: tanto
$G$ quanto a constante $I$ resolvem o sistema linear $Y' = YA +
A^{\mathsf T}Y$ com o mesmo valor inicial, e a unicidade de
Cauchy--Lipschitz para sistemas lineares, \cref{ch:b2:diffeq}, força $G
\equiv I$.)

\emph{Relevância:} o sistema de Frenet $(T, N, B)' = (T, N, B)\,A(s)$
tem a matriz de coeficientes antissimétrica
\[
A = \begin{pmatrix} 0 & -\kappa & 0\\ \kappa & 0 & -\tau\\
0 & \tau & 0\end{pmatrix}
\]
(colunas exprimindo $T', N', B'$). O cálculo acima mostra
que um referencial solução que começa ortonormal \emph{permanece}
ortonormal --- o passo-chave do teorema fundamental que
reconstrói uma curva a partir de $(\kappa, \tau)$.
\end{solution}

\begin{solution}{exo:b2:curves:8}
Pelo~\cref{thm:b2:curves:fundamental} (parte de unicidade), existe uma
função ângulo $\varphi$ de classe $\mathcal{C}^1$ com $T(s) =
(\cos\varphi(s), \sin\varphi(s))$ e $\varphi' = \kappa$. Portanto
\[
\int_0^L \kappa(s)\,\dd s = \varphi(L) - \varphi(0) .
\]
Como o arco é fechado de período $L$, $T(L) = T(0)$:
$(\cos\varphi(L), \sin\varphi(L)) = (\cos\varphi(0),
\sin\varphi(0))$, logo $\varphi(L) - \varphi(0) \in 2\pi\Z$. A
\omterm{thm:b2:curves:frenet2d}{curvatura} total de uma curva fechada é portanto sempre um múltiplo
inteiro de $2\pi$ --- sendo o inteiro o \emph{número de voltas}
da tangente (a quantidade de voltas \omterm{def:b2:metric:complete}{completas} que $T$ dá). Para o
círculo de raio $R$: $\kappa = 1/R$ e $L = 2\pi R$, \omterm{thm:b2:curves:frenet2d}{curvatura}
total $2\pi$, número de voltas $1$; o teorema das tangentes
girantes afirma que esse valor vale para toda curva fechada simples.
\end{solution}

\begin{solution}{exo:b2:curves:9}
Como $\tau \equiv 0$, a curva está num plano
(\cref{prop:b2:curves:torsion}); trabalhe nesse plano. Considere o
centro candidato
\[
c(s) = \gamma(s) + \frac{1}{\kappa}N(s)
\qquad (\kappa \text{ constante}).
\]
Derivando com as fórmulas de Frenet ($N' = -\kappa T + \tau B
= -\kappa T$ aqui):
\[
c'(s) = T + \frac1\kappa(-\kappa T) = 0 ,
\]
logo $c$ é um ponto constante $\Omega$. Então $\norm{\gamma(s) -
\Omega} = \norm{-\frac1\kappa N(s)} = \frac1\kappa$ para todo $s$:
a curva está no círculo de centro $\Omega$ e raio
$1/\kappa$ (em seu plano) e, sendo um arco não constante dele,
é um arco desse círculo.
\end{solution}

\begin{solution}{exo:b2:curves:10}
Para o gráfico $\gamma(x) = (x, x^2/2)$, $\norm{\gamma'(x)} =
\sqrt{1 + x^2}$, logo $L = \int_0^a\sqrt{1+x^2}\,\dd x$.
Substituindo $x = \sinh u$ ($\dd x = \cosh u\,\dd u$, $u$ indo
de $0$ a $u_a = \ln(a + \sqrt{1+a^2})$):
\[
L = \int_0^{u_a}\cosh^2 u\,\dd u
= \frac12\bigl[u + \sinh u\cosh u\bigr]_0^{u_a}
= \frac12\Bigl(\ln\bigl(a + \sqrt{1+a^2}\bigr) +
a\sqrt{1+a^2}\Bigr),
\]
usando $\cosh^2 u = \frac{1 + \cosh 2u}2$ e $\sinh u_a = a$,
$\cosh u_a = \sqrt{1 + a^2}$.
\end{solution}

\begin{solution}{exo:b2:curves:11}
Parametrize por \omterm{def:b2:curves:length}{comprimento de arco}
(\cref{thm:b2:curves:arclength}) e ponha $\lambda(s) =
\langle P - \gamma(s), T(s)\rangle$, uma função
$\mathcal C^1$; como $P$ está sobre a \omterm{def:b2:curves:arc}{reta tangente} em $\gamma(s)$,
o vetor $P - \gamma(s)$ é colinear com $T(s)$, logo $P =
\gamma(s) + \lambda(s)T(s)$. Derivando,
\[
0 = T(s) + \lambda'(s)T(s) + \lambda(s)T'(s)
= \bigl(1 + \lambda'(s)\bigr)T(s) + \lambda(s)T'(s),
\]
e $T'(s) \perp T(s)$ (derive $\norm T^2 = 1$), de modo que
ambas as componentes se anulam: $\lambda' = -1$ e $\lambda T' = 0$.
Ora, $\lambda(s) = c - s$ se anula no máximo uma vez, logo $T' = 0$
num conjunto denso e, por \omterm{def:b2:metric:continuity}{continuidade}, em toda parte: $T$ é um
vetor unitário constante e $\gamma(s) = \gamma(s_0) + (s -
s_0)T$: uma reta (passando por $P$, como deve ser).
\end{solution}

\begin{solution}{exo:b2:curves:12}
(a) A matriz de Frenet
\[
A(s) = \begin{pmatrix} 0 & -\kappa & 0\\ \kappa & 0 & -\tau\\
0 & \tau & 0\end{pmatrix}
\]
tem entradas \omterm{def:b2:metric:continuity}{contínuas}, logo o sistema linear $F' = FA(s)$,
$F(s_0) = F_0$ (uma matriz ortogonal direta) tem solução
única em todo $J$ (\cref{thm:b2:diffeq:linear}). Pelo
\cref{exo:b2:curves:7}, $F(s)$ é ortogonal para todo $s$;
$\det F$ é \omterm{def:b2:metric:continuity}{contínua} com valores em $\{\pm1\}$ e vale
$1$ em $s_0$, logo $F(s)$ é direta para todo $s$.

(b) Leia as linhas $T, N, B$ de $F$ (de modo que $T' = \kappa
N$, $N' = -\kappa T + \tau B$, $B' = -\tau N$) e ponha
$\gamma(s) = \gamma_0 + \int_{s_0}^s T(u)\,\dd u$. Então
$\gamma' = T$ é um vetor unitário: velocidade unitária; $T' = \kappa N$
com $\kappa > 0$ e $N$ unitário ortogonal a $T$, logo $\gamma$
é birregular com \omterm{thm:b2:curves:frenet2d}{curvatura} $\norm{T'} = \kappa$ e
normal principal $N$; a binormal é $T \wedge N = B$
(referencial ortonormal direto), e $B' = -\tau N$ identifica a
\omterm{thm:b2:curves:frenet3d}{torção} como $\tau$.

(c) Sejam $\gamma_1, \gamma_2$ curvas birregulares de velocidade unitária
com os mesmos $(\kappa, \tau)$. Existe uma única isometria
direta $\Phi = \rho + w$ ($\rho \in SO(3)$) levando
$\gamma_1(s_0)$ em $\gamma_2(s_0)$ e o referencial de Frenet de
$\gamma_1$ em $s_0$ no de $\gamma_2$ em $s_0$. A curva
$\Phi\circ\gamma_1$ tem velocidade unitária e os mesmos invariantes
(seu referencial é $\rho$ aplicada ao de $\gamma_1$, e
$\rho$ preserva produtos vetoriais, por ser direta). Ora, os
referenciais de $\Phi\circ\gamma_1$ e de $\gamma_2$ resolvem ambos $F' =
FA(s)$ com o mesmo valor inicial, logo coincidem por
unicidade; em particular as tangentes coincidem e, integrando
a partir do ponto comum $s_0$: $\Phi\circ\gamma_1 = \gamma_2$.
\end{solution}

\begin{solution}{pb:b2:curves:1}
\textbf{1.} O sistema é linear em $(x, y)$ com \omterm{def:b2:linalg:det}{determinante}
$\Delta(t) = a b' - a'b \neq 0$: a regra de Cramer dá a
solução única
\[
x = \frac{c b' - c' b}{\Delta}, \qquad
y = \frac{a c' - a' c}{\Delta}.
\]

\textbf{2.} Como $a(t)x(t) + b(t)y(t) = c(t)$ identicamente,
derivar dá $a'x + b'y + ax' + by' = c'$; a
segunda equação característica mata $a'x + b'y - c'$, logo
$a(t)x'(t) + b(t)y'(t) = 0$: $E'(t)$ é ortogonal a $(a,
b)$, portanto paralelo a $(-b, a)$, a direção de $D_t$. Como
$E(t) \in D_t$ (primeira equação) e $E'(t) \neq 0$, a reta
$D_t$ é exatamente a \omterm{def:b2:curves:arc}{reta tangente} da curva $E$ em $E(t)$.

\textbf{3.} Aqui $(a, b, c) = (t, -1, t^2/2)$, logo $\Delta =
t\cdot0 - 1\cdot(-1) = 1$ e
\[
x = \frac{c b' - c' b}{\Delta} = \tfrac{t^2}2\cdot 0 +
t = t, \qquad
y = \frac{a c' - a' c}{\Delta} = t\cdot t - \tfrac{t^2}2 =
\tfrac{t^2}2 :
\]
a \omterm{pb:b2:curves:1}{envoltória} das \omterm{def:b2:curves:arc}{retas tangentes} da parábola é a
parábola, como devia ser.

\textbf{4.} A reta normal é $\langle M, T(s)\rangle =
\langle\gamma(s), T(s)\rangle$: coeficientes $a = T_1$, $b =
T_2$, $c = \langle\gamma, T\rangle$. Derivando com
Frenet ($T' = \kappa N$): $a' = \kappa N_1$, $b' = \kappa
N_2$ e $c' = \langle T, T\rangle + \langle\gamma, \kappa
N\rangle = 1 + \kappa\langle\gamma, N\rangle$. A segunda
equação característica $\kappa\langle M, N\rangle = 1 +
\kappa\langle\gamma, N\rangle$ dá $\kappa\langle M -
\gamma, N\rangle = 1$. A primeira diz $M - \gamma \perp T$,
logo $M - \gamma = \mu N$ com $\mu = 1/\kappa$: o
ponto característico é $\gamma + \frac1\kappa N$, o \omterm{def:b2:curves:curvature}{centro
de curvatura}, e a \omterm{pb:b2:curves:1}{envoltória} das normais é a evoluta
de \cref{exo:b2:curves:6}. Por fim, $\Delta = T_1\kappa N_2 -
\kappa N_1 T_2 = \kappa\det(T, N) = \kappa \neq 0$.

\textbf{5.} Feixe: o sistema $x\cos\theta + y\sin\theta =
0$, $-x\sin\theta + y\cos\theta = 0$ tem \omterm{def:b2:linalg:det}{determinante} $1$ e
solução $(0,0)$ para todo $\theta$: $E \equiv (0,0)$, $E'
\equiv 0$, e não há curva alguma --- apenas o ponto comum
a todas as retas. Família paralela: $a' = b' = 0$ dá
$\Delta \equiv 0$ e a segunda equação $0 = c'(t)$, que
falha assim que a família de fato se move: nenhum ponto
característico e, de fato, uma família de retas paralelas não toca
curva alguma ao longo de todos os seus membros.

\textbf{6.} A reta por $(\cos t, 0)$ e $(0, \sin t)$
é $\frac x{\cos t} + \frac y{\sin t} = 1$, isto é, $x\sin t +
y\cos t = \sin t\cos t$. Com $(a, b, c) = (\sin t, \cos t,
\sin t\cos t)$: $a' = \cos t$, $b' = -\sin t$, $c' = \cos
2t$, $\Delta = -\sin^2 t - \cos^2 t = -1$. Cramer:
\begin{align*}
x &= \frac{cb' - c'b}{-1} = \sin^2 t\cos t + \cos 2t\cos t
= \cos t\,(\sin^2 t + \cos^2 t - \sin^2 t) = \cos^3 t,\\
y &= \frac{ac' - a'c}{-1} = \sin t\cos^2 t - \sin t\cos 2t
= \sin t\,(\cos^2 t - \cos^2 t + \sin^2 t) = \sin^3 t .
\end{align*}
E $(\cos^3t)^{2/3} + (\sin^3t)^{2/3} = 1$: o \omterm{pb:b2:curves:1}{astroide}.

\textbf{7.} $E'(t) = 3(-\cos^2 t\sin t,\ \sin^2 t\cos t)$
se anula em $t = 0$. Ali, $E''(0) = (-3, 0) \neq 0$ dá
$p = 2$; a componente em $x$ de $E$ é par em $t$, logo
$E'''(0) = (0, 6)$, não colinear: $q = 3$. Par--ímpar: cúspide
de primeira espécie em $(1, 0)$
(\cref{prop:b2:curves:local}), com tangente ao longo do
eixo $x$. As simetrias $x \mapsto -x$, $y \mapsto -y$, $(x,
y) \mapsto (y, x)$ do \omterm{pb:b2:curves:1}{astroide} transportam a cúspide para $(-1,
0)$ e $(0, \pm1)$.

\textbf{8.} $E'(t) = 3\sin t\cos t\,(-\cos t, \sin t)$, logo
$\norm{E'(t)} = 3\abs{\sin t\cos t} = \tfrac32\abs{\sin 2t}$.
Um quadrante: $\int_0^{\pi/2}\tfrac32\sin 2t\,\dd t =
\tfrac32$, e por simetria o \omterm{def:b2:curves:length}{comprimento} total é $4 \cdot
\tfrac32 = 6$.

\textbf{9.} $E(t) - P_t = (\cos^3 t - \cos t,\ \sin^3 t) =
\sin^2 t\,(-\cos t,\ \sin t) = \sin^2 t\,(Q_t - P_t)$. Assim o
ponto de contato é o \omterm{def:b2:affine:barycenter}{baricentro} de $(P_t, \cos^2 t)$ e
$(Q_t, \sin^2 t)$: quando $t$ vai de $0$ a $\pi/2$, ele desliza
da ponta no chão à ponta na parede da escada.

\textbf{10.} No primeiro quadrante a região sob o
\omterm{pb:b2:curves:1}{astroide} tem área $\int_0^1 y\,\dd x$ com $x = \cos^3 t$
decrescendo de $1$ a $0$ quando $t$ vai de $0$ a $\pi/2$:
\[
\int_0^1 y\,\dd x
= \int_{\pi/2}^{0}\sin^3 t\,(-3\cos^2 t\sin t)\,\dd t
= 3\int_0^{\pi/2}\sin^4 t\cos^2 t\,\dd t .
\]
Linearize: $\sin^4 t\cos^2 t = (\sin t\cos t)^2\sin^2 t =
\tfrac18\bigl(\sin^2 2t - \sin^2 2t\cos 2t\bigr)$, e
$\int_0^{\pi/2}\sin^2 2t\,\dd t = \tfrac\pi4$ enquanto
$\int_0^{\pi/2}\sin^2 2t\cos 2t\,\dd t =
\bigl[\tfrac{\sin^3 2t}6\bigr]_0^{\pi/2} = 0$. Logo a
integral vale $\tfrac\pi{32}$, a área do quadrante
$\tfrac{3\pi}{32}$ e a área encerrada $4 \cdot
\tfrac{3\pi}{32} = \tfrac{3\pi}8$.

\textbf{11.} A tangente em $\gamma(t) = (t, t^2/2)$ é
dirigida por $(1, t)$, logo a reta normal é $\{(x, y) : (x -
t) + t\,(y - t^2/2) = 0\}$, isto é, $x + t\,y = t +
\tfrac{t^3}2$.

\textbf{12.} $(a, b, c) = (1, t, t + t^3/2)$: $a' = 0$, $b' =
1$, $c' = 1 + \tfrac32 t^2$, $\Delta = 1$. Então $y = ac' -
a'c = 1 + \tfrac32t^2$ e
\[
x = cb' - c'b = t + \tfrac{t^3}2 - t\Bigl(1 +
\tfrac32t^2\Bigr) = -t^3 .
\]
Eliminando $t$: $t^2 = \tfrac23(y - 1)$ e $x^2 = t^6 =
\tfrac8{27}(y - 1)^3$: uma parábola semicúbica de vértice
$(0, 1)$.

\textbf{13.} $T = (1, t)/\sqrt{1+t^2}$, $N = (-t,
1)/\sqrt{1+t^2}$ e $\kappa = (1+t^2)^{-3/2}$, logo
\[
\gamma + \frac1\kappa N = (t, \tfrac{t^2}2) +
(1+t^2)\,(-t, 1) = \Bigl(-t^3,\ 1 + \tfrac32t^2\Bigr),
\]
a mesma curva: a \omterm{pb:b2:curves:1}{envoltória} das normais é o lugar dos
\omterm{def:b2:curves:curvature}{centros de curvatura}, como prometia a questão 4.

\textbf{14.} $E'(t) = (-3t^2, 3t)$ se anula em $t = 0$;
$E''(0) = (0, 3) \neq 0$ dá $p = 2$ e $E'''(0) = (-6,
0)$ dá $q = 3$: uma cúspide de primeira espécie em $(0, 1)$. O
vértice é onde $\kappa = (1+t^2)^{-3/2}$ é máxima, logo
$\kappa'(0) = 0$ e a velocidade da evoluta
$-\frac{\kappa'}{\kappa^2}N$ se anula exatamente ali: as cúspides
da evoluta ficam nos extremos da \omterm{thm:b2:curves:frenet2d}{curvatura}.

\textbf{15.} A normal de parâmetro $t$ passa por
$(x_0, y_0)$ se e somente se $x_0 + t\,y_0 = t + \tfrac{t^3}2$, isto é,
\[
\frac{t^3}2 + (1 - y_0)\,t - x_0 = 0 ,
\]
uma cúbica em $t$: uma ou três raízes reais (contadas sem
multiplicidade, para pontos genéricos). Sobre o eixo $x_0 = 0$ ela
se fatora como $t\bigl(\tfrac{t^2}2 + 1 - y_0\bigr) = 0$: a
raiz $t = 0$ (o eixo é a normal no vértice), mais
$t = \pm\sqrt{2(y_0 - 1)}$ quando $y_0 > 1$. Assim: uma normal
para $y_0 < 1$, três para $y_0 > 1$ e, em $y_0 = 1$, a
raiz tripla marca a cúspide da evoluta. Em geral, uma
raiz dupla da cúbica significa que o ponto satisfaz tanto a
equação da reta quanto sua derivada em $t$ --- ele está \emph{sobre a
\omterm{pb:b2:curves:1}{envoltória}}: a evoluta é precisamente a fronteira entre as
regiões de uma e de três normais.

\textbf{16.} A reflexão no espelho inverte a componente normal
da direção e conserva a tangencial:
escrevendo $u = \langle u, n\rangle n + u_{\mathrm{tan}}$, a
direção refletida é $u_{\mathrm{tan}} - \langle u,
n\rangle n = u - 2\langle u, n\rangle n$. Aqui $\langle u,
n\rangle = \cos\theta$, logo
\[
v = (1, 0) - 2\cos\theta\,(\cos\theta, \sin\theta)
= (1 - 2\cos^2\theta,\ -2\sin\theta\cos\theta)
= -(\cos2\theta,\ \sin2\theta).
\]

\textbf{17.} O raio refletido passa por $P_\theta =
(\cos\theta, \sin\theta)$ com direção $(\cos2\theta,
\sin2\theta)$; um vetor normal é $(-\sin2\theta,
\cos2\theta)$, de modo que a reta é
\[
-\sin2\theta\,(x - \cos\theta) + \cos2\theta\,(y -
\sin\theta) = 0,
\]
e a constante vale $-\sin2\theta\cos\theta +
\cos2\theta\sin\theta = -\sin\theta$: multiplicando por $-1$,
$x\sin2\theta - y\cos2\theta = \sin\theta$.

\textbf{18.} $(a, b, c) = (\sin2\theta, -\cos2\theta,
\sin\theta)$: $a' = 2\cos2\theta$, $b' = 2\sin2\theta$, $c' =
\cos\theta$, $\Delta = 2\sin^22\theta + 2\cos^22\theta = 2$.
Cramer e, em seguida, as fórmulas de produto em soma:
\begin{align*}
x &= \frac{2\sin\theta\sin2\theta +
\cos\theta\cos2\theta}{2}
= \frac{(\cos\theta - \cos3\theta) + \frac12(\cos\theta +
\cos3\theta)}{2} = \frac{3\cos\theta - \cos3\theta}4,\\
y &= \frac{\sin2\theta\cos\theta - 2\cos2\theta\sin\theta}2
= \frac{\frac12(\sin3\theta + \sin\theta) - (\sin3\theta -
\sin\theta)}2 = \frac{3\sin\theta - \sin3\theta}4 :
\end{align*}
a \omterm{pb:b2:curves:1}{nefroide}, uma curva fechada com duas cúspides.

\textbf{19.} Derivando e fatorando com $\sin3\theta
- \sin\theta = 2\cos2\theta\sin\theta$, $\cos\theta -
\cos3\theta = 2\sin2\theta\sin\theta$:
\[
E'(\theta) = \tfrac34\bigl(\sin3\theta - \sin\theta,\
\cos\theta - \cos3\theta\bigr)
= \tfrac32\sin\theta\,(\cos2\theta, \sin2\theta),
\]
paralelo à direção refletida da questão 16: cada
raio refletido é tangente à cáustica, como exige a propriedade
de \omterm{pb:b2:curves:1}{envoltória}. $E' = 0$ exatamente em $\sin\theta = 0$:
$E(0) = (\tfrac12, 0)$ e $E(\pi) = (-\tfrac12, 0)$, as duas
cúspides. Pondo $y = 0$ na equação da reta: $x\sin2\theta =
\sin\theta$, logo $x = \frac1{2\cos\theta} \to \frac12$ quando
$\theta \to 0$: os raios paraxiais focalizam à distância $R/2$ do
centro --- a distância focal do espelho esférico.

\textbf{20.} $\norm{E'(\theta)} = \tfrac32\abs{\sin\theta}$,
logo o \omterm{def:b2:curves:length}{comprimento} vale $\tfrac32\int_0^{2\pi}\abs{\sin\theta}\,
\dd\theta = \tfrac32\cdot4 = 6$. Perto de $\theta = 0$, com
$\cos k\theta = 1 - \tfrac{k^2\theta^2}2 + O(\theta^4)$ e
$\sin k\theta = k\theta - \tfrac{k^3\theta^3}6 +
O(\theta^5)$:
\[
x - \tfrac12 = \frac{3\theta^2 + O(\theta^4)}{4}
= \tfrac34\theta^2 + O(\theta^4),
\qquad
y = \frac{4\theta^3 + O(\theta^5)}{4}
= \theta^3 + O(\theta^5) :
\]
$p = 2$, $q = 3$, uma cúspide de primeira espécie apontando ao longo do
eixo --- o ponto brilhante da cáustica da xícara de café.

\textbf{21.} Parametrize os pontos iluminados por
$\Phi(\theta, r) = P_\theta + r\,v_\theta$ (posição ao longo
de cada raio refletido). O \omterm{def:b2:linalg:det}{determinante} jacobiano
$\det(\partial_\theta\Phi, \partial_r\Phi) =
\det(P_\theta' + rv_\theta', v_\theta)$ é \omterm{def:b2:affine:subspace}{afim} em $r$ e
se anula para exatamente um $r = r_*(\theta)$ --- e
$\Phi(\theta, r_*(\theta))$ é o ponto característico, pois
ali a direção do raio e a variação da família se tornam
dependentes. Fora da \omterm{pb:b2:curves:1}{envoltória} a aplicação é um difeomorfismo local
e um ponto logo dentro da cáustica é atingido por dois raios
próximos (duas soluções $\theta$), um ponto de fora por nenhum daquela
parte da família; sobre a cáustica os dois se fundem. A
intensidade luminosa é inversamente proporcional ao valor absoluto
do jacobiano, de modo que ela explode ao longo da \omterm{pb:b2:curves:1}{envoltória}: a cáustica é a
curva brilhante, mais brilhante de todas na cúspide, onde a
degenerescência é máxima.

\textbf{22.} $x' = 1 - \cos t$, $y' = \sin t$, $x'' = \sin
t$, $y'' = \cos t$, logo $x'y'' - y'x'' = \cos t - 1$ e
$\norm{\gamma'}^2 = 2(1 - \cos t) = 4\sin^2\tfrac t2$;
pela~\cref{prop:b2:curves:kappaformula},
\[
\kappa(t) = \frac{-(1 - \cos t)}{8\sin^3\tfrac t2}
= \frac{-2\sin^2\tfrac t2}{8\sin^3\tfrac t2}
= -\frac1{4\sin\tfrac t2}
\qquad (0 < t < 2\pi).
\]
Com $T = (\sin\tfrac t2, \cos\tfrac t2)$ (divida $\gamma'$
por $2\sin\tfrac t2$) e $N = (-\cos\tfrac t2, \sin\tfrac
t2)$:
\[
\gamma + \frac1\kappa N
= \gamma - 4\sin\tfrac t2\,\Bigl(-\cos\tfrac t2,\
\sin\tfrac t2\Bigr)
= \bigl(t - \sin t + 2\sin t,\ 1 - \cos t - 2(1 - \cos
t)\bigr),
\]
isto é, $c(t) = (t + \sin t,\ \cos t - 1)$. Substituindo $t = u
+ \pi$:
\[
c = \bigl(u + \pi - \sin u,\ -\cos u - 1\bigr)
= \bigl((u - \sin u) + \pi,\ (1 - \cos u) - 2\bigr) :
\]
a cicloide $\gamma(u)$ transladada por $(\pi, -2)$. A
evoluta de uma cicloide é uma cicloide congruente, pendurada um
nível abaixo.

\textbf{23.} $R(t) = 1/\abs{\kappa(t)} = 4\sin\tfrac t2$, logo
$R(\pi) = 4$: metade do \omterm{def:b2:curves:length}{comprimento} $8$ do arco, calculado no
\cref{exo:b2:curves:1}. A velocidade da evoluta $c'(t) = (1 +
\cos t, -\sin t)$ se anula em $t = \pi$: a cúspide é $c(\pi)
= (\pi, -2)$, diretamente abaixo do ápice $\gamma(\pi) = (\pi,
2)$, à distância $4 = R(\pi)$, exatamente o \omterm{def:b2:curves:length}{comprimento} do
raio osculador ali.

\textbf{24.} De \cref{exo:b2:curves:6}, $c'(s) =
-\frac{\kappa'(s)}{\kappa(s)^2}N(s) = R'(s)\,N(s)$, logo
$\norm{c'(s)} = \abs{R'(s)}$ e, para $R$ monótona,
\[
\int_{s_0}^{s_1}\norm{c'(s)}\,\dd s =
\Bigl|\int_{s_0}^{s_1}R'(s)\,\dd s\Bigr| = \abs{R(s_1) -
R(s_0)} .
\]
Digamos que $R$ decresça. Um fio estendido ao longo da evoluta além de
$c(s_0)$ e prolongado pelo segmento de $c(s_0)$ a
$\gamma(s_0)$ (que é tangente à evoluta, pela questão 4)
tem, ao ser desenrolado até $c(s)$ e esticado, parte
reta de \omterm{def:b2:curves:length}{comprimento} $R(s_0) - \bigl(R(s_0) - R(s)\bigr) = R(s)$
apontando de $c(s)$ ao longo da normal --- caindo exatamente sobre
$\gamma(s)$: a ponta livre traça a curva original
(``evolvente''). Huygens pendurou um pêndulo entre duas faces
cicloidais: o cordão se enrola na evoluta, de modo que o peso descreve uma
cicloide --- a tautócrona, cujo período de oscilação não
depende da amplitude.

\textbf{25.} Tangentes da parábola: $\Delta = 1$ e $E' =
(1, t) \neq 0$ em toda parte --- \omterm{pb:b2:curves:1}{envoltória} suave (a própria
parábola). Escada que escorrega: $\Delta = -1$, mas $E' =
\tfrac32\sin 2t\,(-\cos t, \sin t)$ se anula nas
extremidades do quadrante --- as quatro cúspides do \omterm{pb:b2:curves:1}{astroide}. Normais da
parábola e da cicloide: $\Delta = \kappa \neq 0$, e
$E' = R'N$ se anula exatamente onde a \omterm{thm:b2:curves:frenet2d}{curvatura} é extremal
--- cúspides das evolutas em $(0,1)$ e $(\pi, -2)$. Cáustica:
$\Delta = 2$, e $E' = \tfrac32\sin\theta\,(\cos2\theta,
\sin2\theta)$ se anula em $\theta = 0, \pi$ --- as duas cúspides
da \omterm{pb:b2:curves:1}{nefroide}, os pontos focais do espelho. Extremos de
\omterm{thm:b2:curves:frenet2d}{curvatura} \emph{têm} de produzir cúspides na \omterm{pb:b2:curves:1}{envoltória} das
normais, pois a velocidade da evoluta é $R'N$: é por isso que
a evoluta da elipse tem quatro cúspides (quatro vértices),
e as famílias degeneradas (feixe, paralelas) são os casos
em que a máquina devolve um ponto ou nada.
\end{solution}
